% ============================================================
\chapter{Applications}
\label{ch:applications}
% ============================================================

Fixed point theorems are not abstract curiosities: they are the tools that, behind the scenes, make a large part of classical analysis work. The Picard--Lindel\"of existence and uniqueness theorem for ordinary differential equations? It is Banach's contraction principle applied to an integral operator. Fredholm integral equations? Banach again. The implicit function theorem, pillar of differential calculus? The contraction principle, once more. This chapter gathers these classical applications, showing how a single abstract result --- the existence of a fixed point for a contraction --- irrigates entire branches of applied mathematics.

\begin{intuition}
This chapter illustrates the power of fixed point theorems by applying them to three
classical areas of analysis: ordinary differential equations (via the Picard--Lindel\"of
theorem as a consequence of Banach's contraction principle), integral equations (Fredholm
and Volterra), and the implicit function theorem. In each case, the problem is reformulated
as a fixed point problem in an appropriate function space.
\end{intuition}

% ------------------------------------------------------------
\section{Ordinary Differential Equations: Picard--Lindel\"of}
% ------------------------------------------------------------

\begin{theorem}[Picard--Lindel\"of (Cauchy--Lipschitz)]
\label{thm:picard-lindelof}
\index{Picard--Lindel\"of theorem}
Let $f\colon [t_0 - a, t_0 + a] \times \overline{B}(y_0, b) \to \R^n$ be a continuous
map, Lipschitz in $y$ with constant $L$:
\[
  \norm{f(t, y_1) - f(t, y_2)} \leq L \norm{y_1 - y_2}
\]
for all $t, y_1, y_2$. Let $M = \sup \norm{f}$ and $\delta = \min(a, b/M)$. Then the
Cauchy problem
\[
  y'(t) = f(t, y(t)), \qquad y(t_0) = y_0
\]
has a unique solution $y \in C^1([t_0 - \delta, t_0 + \delta], \R^n)$.
\end{theorem}

\begin{proof}
The Cauchy problem is equivalent to the integral equation
\[
  y(t) = y_0 + \int_{t_0}^{t} f(s, y(s)) \, ds =: (Ty)(t).
\]
Consider the space $E = \{y \in C([t_0 - \delta, t_0 + \delta], \R^n) :
\norm{y(t) - y_0} \leq b \; \forall t\}$ equipped with the norm
\[
  \norm{y}_\lambda = \sup_{t} e^{-\lambda \abs{t - t_0}} \norm{y(t) - y_0}
\]
for $\lambda > L$ to be chosen.

\emph{$T$ maps $E$ into $E$}: For $y \in E$,
$\norm{(Ty)(t) - y_0} \leq M \abs{t - t_0} \leq M\delta \leq b$.

\emph{$T$ is a contraction}: For $y_1, y_2 \in E$,
\begin{align*}
  e^{-\lambda\abs{t-t_0}} \norm{(Ty_1)(t) - (Ty_2)(t)}
  &\leq e^{-\lambda\abs{t-t_0}} \int_{t_0}^{t} L \norm{y_1(s) - y_2(s)} \, ds \\
  &\leq L \norm{y_1 - y_2}_\lambda \int_{t_0}^{t} e^{\lambda(\abs{s-t_0} - \abs{t-t_0})} ds \\
  &\leq \frac{L}{\lambda} \norm{y_1 - y_2}_\lambda.
\end{align*}
Choosing $\lambda > L$, we get $\norm{Ty_1 - Ty_2}_\lambda \leq
\frac{L}{\lambda} \norm{y_1 - y_2}_\lambda$ with $L/\lambda < 1$.

By Banach's contraction principle, $T$ has a unique fixed point in $E$.
\end{proof}

\begin{remark}
The norm $\norm{\cdot}_\lambda$ (Bielecki norm) is a classical trick that avoids having to
restrict the time interval to obtain contraction.
\end{remark}

\begin{corollary}[Continuous dependence]
The solution $y(t; t_0, y_0)$ depends continuously on the initial data $(t_0, y_0)$.
\end{corollary}

% ------------------------------------------------------------
\section{Fredholm Integral Equations}
% ------------------------------------------------------------

\begin{definition}[Fredholm equation of the second kind]
\index{Fredholm equation}
Let $K \in C([a,b]^2, \R)$ be a \emph{kernel}, $f \in C([a,b], \R)$, and
$\lambda \in \R$. The \emph{Fredholm equation of the second kind} is:
\[
  u(t) = f(t) + \lambda \int_a^b K(t,s) \, u(s) \, ds.
\]
\end{definition}

\begin{theorem}[Existence and uniqueness for Fredholm]
\label{thm:fredholm}
If $\abs{\lambda} < \frac{1}{(b-a) \norm{K}_\infty}$, then the Fredholm equation has a
unique solution $u \in C([a,b], \R)$.
\end{theorem}

\begin{proof}
Define the operator $T\colon C([a,b]) \to C([a,b])$ by
\[
  (Tu)(t) = f(t) + \lambda \int_a^b K(t,s) \, u(s) \, ds.
\]
Then:
\[
  \norm{Tu_1 - Tu_2}_\infty \leq \abs{\lambda} (b-a) \norm{K}_\infty \norm{u_1 - u_2}_\infty.
\]
Under the hypothesis $\abs{\lambda} (b-a) \norm{K}_\infty < 1$, the operator $T$ is a
contraction on the Banach space $C([a,b])$, and Banach's theorem applies.
\end{proof}

\begin{example}
The equation $u(t) = \sin(t) + \frac{1}{4} \int_0^1 ts \, u(s) \, ds$ has kernel
$K(t,s) = ts$ with $\norm{K}_\infty = 1$. Since
$\abs{\lambda}(b-a)\norm{K}_\infty = \frac{1}{4} < 1$, the solution exists and is unique.
\end{example}

% ------------------------------------------------------------
\section{Volterra Integral Equations}
% ------------------------------------------------------------

\begin{definition}[Volterra equation of the second kind]
\index{Volterra equation}
The \emph{Volterra equation of the second kind} is:
\[
  u(t) = f(t) + \lambda \int_a^t K(t,s) \, u(s) \, ds.
\]
The difference from Fredholm is that the upper limit of integration is $t$ (not $b$).
\end{definition}

\begin{theorem}[Existence and uniqueness for Volterra]
\label{thm:volterra}
For any $\lambda \in \R$ and any continuous kernel $K$, the Volterra equation of the second
kind has a unique solution $u \in C([a,b], \R)$.
\end{theorem}

\begin{proof}
The operator $T$ defined by $(Tu)(t) = f(t) + \lambda \int_a^t K(t,s) u(s) ds$ is not
necessarily a contraction for the uniform norm. However, we use the Bielecki norm:
\[
  \norm{u}_\mu = \sup_{t \in [a,b]} e^{-\mu(t-a)} \abs{u(t)}.
\]
Then:
\begin{align*}
  e^{-\mu(t-a)} \abs{(Tu_1)(t) - (Tu_2)(t)}
  &\leq \abs{\lambda} \norm{K}_\infty \int_a^t e^{-\mu(t-s)} e^{-\mu(s-a)} \abs{u_1(s) - u_2(s)} \, ds \\
  &\leq \abs{\lambda} \norm{K}_\infty \norm{u_1 - u_2}_\mu \int_a^t e^{-\mu(t-s)} ds \\
  &\leq \frac{\abs{\lambda} \norm{K}_\infty}{\mu} \norm{u_1 - u_2}_\mu.
\end{align*}
Choosing $\mu > \abs{\lambda} \norm{K}_\infty$, $T$ is a contraction.
\end{proof}

\begin{attention}
Unlike Fredholm, the Volterra equation always has a unique solution, with no restriction on
$\lambda$. This is a consequence of the ``triangular'' character of the Volterra operator.
\end{attention}

% ------------------------------------------------------------
\section{The Implicit Function Theorem}
% ------------------------------------------------------------

\begin{theorem}[Implicit function theorem via fixed point]
\label{thm:implicit}
\index{implicit function theorem}
Let $F\colon U \to \R^m$ be a $C^1$ map where $U$ is an open subset of $\R^n \times \R^m$.
Let $(a, b) \in U$ with $F(a, b) = 0$ and $D_y F(a, b)$ invertible. Then there exist
neighborhoods $V$ of $a$ and $W$ of $b$ and a unique $C^1$ map $\varphi\colon V \to W$
such that
\[
  F(x, \varphi(x)) = 0 \quad \forall x \in V.
\]
\end{theorem}

\begin{proof}
We reduce to a fixed point problem. Set $A = D_y F(a,b)$ and
$G(x,y) = y - A^{-1} F(x,y)$. Then $F(x,y) = 0$ is equivalent to $y = G(x,y)$.

We show that for $x$ near $a$, $y \mapsto G(x,y)$ is a contraction near $b$. We have
$D_y G(a,b) = I - A^{-1} D_y F(a,b) = 0$. By continuity of $D_y G$, there exists $r > 0$
such that $\norm{D_y G(x,y)} \leq 1/2$ in a neighborhood of $(a,b)$.

Thus, for fixed $x$ near $a$, $y \mapsto G(x,y)$ is a contraction with constant $1/2$ on
$\overline{B}(b, r)$ (adjusting $r$ if necessary). By Banach's theorem, there exists a
unique $y = \varphi(x)$ with $G(x, \varphi(x)) = \varphi(x)$, i.e.,
$F(x, \varphi(x)) = 0$.

The $C^1$ regularity of $\varphi$ is obtained by differentiating $F(x, \varphi(x)) = 0$,
giving $D\varphi(x) = -[D_y F(x, \varphi(x))]^{-1} D_x F(x, \varphi(x))$.
\end{proof}

\begin{remark}
This fixed point proof is more constructive than the classical proof (via the inverse
function theorem) because it provides an approximation algorithm: the sequence
$y_{n+1} = G(x, y_n)$ converges to $\varphi(x)$.
\end{remark}

% ------------------------------------------------------------
\section{Complements: Newton's Method and Fixed Points}
% ------------------------------------------------------------

\begin{proposition}[Newton's method as fixed point iteration]
Newton's method for solving $F(x) = 0$ in $\R^n$ reads:
\[
  x_{n+1} = x_n - [DF(x_n)]^{-1} F(x_n) =: T(x_n).
\]
Under regularity conditions (Lipschitz $DF$, invertibility of $DF(x^*)$), $T$ is not a
global contraction, but it is locally contractive near the solution $x^*$, with a
quadratic convergence rate: $\norm{x_{n+1} - x^*} \leq C \norm{x_n - x^*}^2$.
\end{proposition}

\begin{keyformulas}[title=\textbf{Summary of applications}]
\begin{center}
\renewcommand{\arraystretch}{1.3}
\begin{tabular}{|l|l|l|}
\hline
\textbf{Problem} & \textbf{Theorem used} & \textbf{Function space} \\
\hline
ODE (Picard--Lindel\"of) & Banach & $C([t_0-\delta, t_0+\delta], \R^n)$ \\
Fredholm & Banach & $C([a,b], \R)$ \\
Volterra & Banach (Bielecki norm) & $C([a,b], \R)$ \\
Implicit functions & Banach & $\R^m$ \\
\hline
\end{tabular}
\end{center}
\end{keyformulas}

% ------------------------------------------------------------
\section{Exercises}
% ------------------------------------------------------------

\begin{exercise}
Consider the ODE $y' = y^2$, $y(0) = 1$. Compute the first three Picard iterates
$y_0(t) = 1$, $y_{n+1}(t) = 1 + \int_0^t y_n(s)^2 ds$. Compare with the exact solution
$y(t) = \frac{1}{1-t}$.
\end{exercise}

\begin{exercise}
Solve the Fredholm equation $u(t) = 1 + \lambda \int_0^1 u(s) ds$ explicitly for all
$\lambda \neq 1$. Verify that the solution diverges as $\lambda \to 1$.
\end{exercise}

\begin{exercise}
Solve the Volterra equation $u(t) = 1 + \int_0^t u(s) ds$ by recognizing an ODE. Verify
the result by the method of successive iterations.
\end{exercise}

\begin{exercise}
Show that the equation $F(x, y) = x^2 + y^2 - 1 = 0$ defines $y$ as a function of $x$
near $(0, 1)$. Compute $\varphi'(0)$ and $\varphi''(0)$.
\end{exercise}

\begin{exercise}
Let $K(t,s) = e^{ts}$ and $f(t) = 1$. Study the convergence of the iterates of the
operator $T$ associated with the Fredholm equation
$u(t) = 1 + \lambda \int_0^1 e^{ts} u(s) ds$ as a function of $\lambda$.
\end{exercise}

\begin{exercise}[Systems of ODEs]
Apply the Picard--Lindel\"of theorem to the system $y_1' = y_2$, $y_2' = -y_1$,
$y_1(0) = 1$, $y_2(0) = 0$. Compute the Picard iterates explicitly and show they converge
to $(\cos t, -\sin t)$.
\end{exercise}
